\documentclass[10pt,a4paper]{amsart}
\usepackage[margin=19mm]{geometry}
\usepackage{amsmath,amssymb,mathtools}
\usepackage[scr=rsfs,cal=euler]{mathalfa}
\usepackage{xfrac}
\usepackage[unicode,hidelinks,bookmarks=false]{hyperref}
\newcommand{\ie}{i.e.\ }
\newcommand{\cf}{cf.\ }
\newcommand{\resp}{respectively\ }
\newcommand{\Subsection}{\subsection}
\providecommand{\nobreakdash}{\nobreak\hbox{-}}
\setlength{\emergencystretch}{2em}
\multlinegap0pt
\usepackage{amssymb}
\usepackage{enumerate}
\newtheorem{lem}{Lemma}
\newtheorem{cor}{Corollary}
\newtheorem{prop}{Proposition}
\newcommand\F{\mathbb{F}}
\newcommand\Z{\mathbb{Z}}
\newcommand\gam{\gamma}
\newcommand\la{\langle}
\newcommand\lla{\la\!\la}
\newcommand\lm{\lambda}
\newcommand\ra{\rangle}
\newcommand\rra{\ra\!\ra}
\title{Groups of generalized Moufang type and $\Z_2$-graded algebras.}
\author{Ilya Gorshkov}
\date{Source: arXiv:2603.01988v1 (CC BY 4.0)}
\begin{document}
\maketitle

\noindent\emph{Source and licence.} Groups of generalized Moufang type and $\Z_2$-graded algebras.. Version arXiv:2603.01988v1 (CC BY 4.0), distributed by arXiv under the Creative Commons Attribution 4.0 International licence, http://creativecommons.org/licenses/by/4.0/, as recorded in the arXiv licence metadata for this exact version and shown on the abstract page https://arxiv.org/abs/2603.01988v1. Full licence notice: \texttt{licenses/arxiv-2603.01988-CC-BY-4.0.txt}.\par\smallskip

\noindent\textit{Editorial scope.} Counted unit: the article's prop twograd (source0.tex lines 901--904). The article's complete original proof of it is delivered with it (source0.tex lines 905--1034, one \texttt{proof} environment of 130 lines). No mathematics is written, repaired or re-derived: the statement and the proof are the article's own lines, in the article's own notation, and no retained line is substituted or re-flowed.  Retained uncounted as the source-local context the proof needs: lines 596--616; lines 744--747; lines 751--754; lines 768--770; lines 884--886.  Nothing was omitted from any retained span, so the package carries no omission record.  No retained line contains a citation, so the wrapper carries no bibliography and no bibliography entry.  No retained line contains a figure environment, an image-inclusion command or drawing code, so no image is delivered and the package declares no asset dependency.  Editorial formatting only, in addition: an AMS \texttt{amsart} preamble that restates the article's own declarations and macros used by the retained lines, namely the environments \texttt{lem}, \texttt{cor}, \texttt{prop} on separate counters, together with the standard packages those lines need.  The retained environments print in this document as Lemma 1, Corollary 1, Proposition 1 in order of appearance, whereas the article prints the same items with the numbering of its own full document.  The complete native source of the article as distributed in the arXiv e-print for this exact version is bundled unchanged as \texttt{sources/195585/source0.tex}.
\begin{lem}\label{prod1}
Products of basis elements of the space $A_a^+$:
%Произведения базисных элементов пространства $A_a^+$.
\begin{enumerate}
\item let $p>5$, If $i\leq \frac{p-3}{4}$ then $y_i(a)*a=(1-\eta)(y_{2i}(a)-y_1(a))$ else $y_i(a)*a=(1-\eta)(y_{p-3-2i}(a)-y_1(a))$;

\item let $p=5$, we have $y_1(a)*a=(\eta-1)y_1(a)$;

\item $z(a)*y_i(a)= \frac{\eta(p-3)+1}{p-2}y_i(a)$;

\item $y_i(a)*z(a)=y_i(a)+y_i(a)*a\in L(y_1,...,y_{(p-3)/2})$;

\item $z(a)*z(a)=\frac{1+2\eta(p-2)}{(p-2)^2}a+(\eta(p-1)^2+p)\gam\in L(a,z)$; 

\item $z(a)*a=a*z(a)=(1+(p-2)\eta)z(a)$;

\item $y_i(a)*y_i(a), y_i(a)*y_j(a)\in A_a^+$;

\item if $p=5$, then $y_1(a)*y_1(a)=(1-4\eta)z(a)-\frac{16}{3}(1-\eta)a.$ 
\end{enumerate}
\end{lem}

\begin{cor}\label{twogenaxial}
	%Разложение алгебры $A$ на $A^+_a$ и $A^-_a$ является $2$-градуировкой алгебры $A$, более того верны следующие утверждения:
	For each idempotent $c\in T$, the decomposition of the algebra $A$ into a direct product of the subspaces $A^+_c$ and $A^-_c$ is a $\Z_2$-grading of the algebra $A$.
\end{cor}

\begin{lem}\label{xxa}
For any $x\in A$, $x+x^a\in A^+_a$.
%Для любого $x\in A$ выполнено $x+x^a\in A^+_a$.
\end{lem}

\begin{lem}\label{Mi2}
	For any $c,d\in T$, we have $\tau_c(d)=d^c$.
\end{lem}

\begin{lem}\label{aaut}
	Let $x,y\in A, a\in T$, then $(x*y)^a=x^a*y^a$.
\end{lem}

\begin{prop}\label{twograd}
 $(A_a^+,A_a^-)$ is a $\Z_2$-grading of $A$. 
 %С каждым идемпотентом $a\in T$ связана нетривиальная $\Z_2$-градуировка $GA_a=(A_a^+,A_a^-)$ алгебры $A_{\F}(G,T,\eta)$. Если $b\in T\setminus\{a\}$, то $(A_a^+,A_a^-)\neq(A_b^+,A_b^-)$. 	    
\end{prop}	

\begin{proof}
	If $G$ is generated by two elements of $T$, then the proposition follows from Corollary \ref{twogenaxial}. 
    Since $A_a^+\oplus A_a^-$ contains the basis $A$, then $A_a^+\oplus A_a^-=A$. We divide the proof of the proposition into several lemmas.
    
	\begin{lem}\label{xxagen}
		$x+x^a\in A^+_a$
	\end{lem}
	\begin{proof}
	Since $x$ is a sum of elements of $T$, it suffices to show that $b+b^a\in A^+_a$ for $b\in T$. This statement is proved in Lemma \ref{xxa}.
\end{proof}
	
	Let $b,c\in T$ and $c\not \in I(a,b)$, 
    $(a,z, y_1,...,y_k, x_1,...,x_{k+1})=BE(L_a, a, b)$, and
    
    $(a, z', y'_1,...,y'_k, x'_1,...,x'_{k+1})=BE(L_a, a, c)$.
		
	\begin{lem}
	$x_i*x'_j\in A^+_a$
	\end{lem}
	\begin{proof}
	
	%Poskol'ku mnozhestva $\{x_1,...,x_{k+1}\}$ i $\{x'_1,...,x'_{k+1}\}$ ne peresekayutsya, to bez ogranicheniya obshchnosti, my mozhem schitat', chto $i=j=1$. Imeyem $$W=bet_1*bet'_1=(-b+b^a)*(-c+c^a)=b*c+ b^a*c^a-b^a*c-b*c^a.$$ Iz lemmy \ref{aaut} sleduyet, chto $b^a*c^a=(b*c)^a$, $b^a*c=(b*c^a)^a$. % $$b*c^a=c^{ab}+\eta\sum_{x\in I(b,c^a)\setminus c^{ab}}x$$ $$W=b*c+(b*c)^a-b*c^a -(b*c^a)^a$$ Pust' $x=b*c-b*c^a$. Sledovatel'no $W=x+x^a$. Iz lemmy \ref{xxagen} sleduyet, chto $W\in A^+$.

	Since the sets $\{x_1,...,x_{k+1}\}$ and $\{x'_1,...,x'_{k+1}\}$ are disjoint, without loss of generality, we can assume that $i=j=1$.
	We have
	$$W=x_1*x'_1=(-b+b^a)*(-c+c^a)=b*c+ b^a*c^a-b^a*c-b*c^a.$$
	From Lemma \ref{aaut} it follows that $b^a*c^a=(b*c)^a$ and $b^a*c=(b*c^a)^a$. We have
	
	% $$b*c^a=c^{ab}+\eta\sum_{x\in I(b,c^a)\setminus c^{ab}}x$$
	
	$$W=b*c+(b*c)^a-b*c^a -(b*c^a)^a$$
	Let $x=b*c-b*c^a$. Therefore, $W=x+x^a$.
	From Lemma \ref{xxagen} it follows that $W\in A^+$.
	
	%Поскольку множества $\{x_1,...,x_{k+1}\}$ и $\{x'_1,...,x'_{k+1}\}$ не пересекаются, то без ограничения общности, мы можем считать, что $i=j=1$. 
	%Имеем 
	%$$W=bet_1*bet'_1=(-b+b^a)*(-c+c^a)=b*c+ b^a*c^a-b^a*c-b*c^a.$$
	%Из леммы \ref{aaut} следует, что $b^a*c^a=(b*c)^a$, $b^a*c=(b*c^a)^a$.
	
%	$$b*c^a=c^{ab}+\eta\sum_{x\in I(b,c^a)\setminus c^{ab}}x$$
		
	%$$W=b*c+(b*c)^a-b*c^a -(b*c^a)^a$$
	%Пусть $x=b*c-b*c^a$. Следовательно $W=x+x^a$.
	%Из леммы \ref{xxagen} следует, что $W\in A^+$.
	
\end{proof}
	
	Thus, the product of two elements of $A_a^-$ lies in $A_a^+$.
	
	Let's verify that the product of an element of $A^-_a$ and an element of $A^+_a$ lies in $A^-_a$. Since $A^-_a$ is a eigenspace with respect to the operator $L_a$, $a*x$ lies in $A^-_a$ for any element $x\in A^-_a$.
	
	\begin{lem}\label{xmx}
		Let $x\in A$. We have $x-x^a\in A^-_a$.
	\end{lem}
	\begin{proof}
		Since $x$ is a linear combination of elements of $T$, to prove the statement of the lemma it suffices to show that $b-b^a\in A^-_a$, where $b\in T$. This statement proved in Lemma \ref{Mi2}.
		%Поскольку $x$ является линейной комбинацией элементов из $T$, для доказательства утверждения достаточно показать, что $b-b^a\in A^-$, где $b\in T$. Это утверждение следует из Леммы \ref{Mi2}.
	\end{proof}
	\begin{lem}\label{betia}
		Let $x\in A^-_a$. We have $x*a\in A^-_a$.
	\end{lem}
	\begin{proof}
The space $A_a^-$ has a basis of eigenvectors $L_a$. We can assume that $x$ is an eigenvector of $L_a$ with eigenvalue $-\lm_2$ lying in the subalgebra $\lla a,b\rra$ for some $b\in T$.
The assertion of the lemma follows from Corollary \ref{twogenaxial}.
%	Пространство $A_a^-$ имеет базис из собственных векторов $L_a$. Можно считать, что $x$ это собственный вектор оператора $L_a$ с собственным числом $-\lm_2$ лежащий в подалгебре $\lla a,b\rra$ для некоторого $b\in T$.
%Утверждение леммы следует из \ref{twogenaxial}.
\end{proof}	
		
		
	
	\begin{lem}\label{cd}
	Let $c,d\in T$. We have $(c+c^a)*(d-d^a)\in A^-_a$.
	\end{lem}
	\begin{proof}
	We have $$W=(c+c^a)*(d-d^a)=a*d-c^a*d^a+c^a*d-c*d^a.$$
From Lemma \ref{aaut} it follows that $c^a*d^a=(c*d)^a$ and $c*d^a=(c^ad)^a$.
From Lemma \ref{xmx} it follows that $W\in A^-_a$.

	%Из леммы \ref{aaut} следует, что $c^a*d^a=(c*d)^a$ и $c*d^a=(c^ad)^a$.
	% Из леммы \ref{xmx} следует, что $W\in A^-$. 	
	\end{proof}
	
	\begin{lem}\label{dc}
	Let $c,d\in T$. We have $(c-c^a)*(d+d^c)\in A^-_a$.
	\end{lem}
	\begin{proof}
The proof is similar to the proof of Lemma \ref{cd}.
%	Доказательство аналогично доказательству леммы \ref{cd}.
	\end{proof}

	\begin{lem}\label{zbet}
	We have $z*x'_i, x'_i*z\in A^-_a$.
	\end{lem}
	\begin{proof}
Note that $z$ is linearly expressible in terms of $a$ and elements of the form $x+x^a$. Thus, the assertion of the lemma follows from Lemmas \ref{cd} and \ref{dc}.

%		Заметим, что $z$ линейно выражается через $a$ и элементы вида $x+x^a$. Таким образом, утверждение леммы следует из лемм \ref{cd} и \ref{dc}. 
	\end{proof}
	
	\begin{lem}\label{alpbet}
	We have $x_i*y'_j, y_i*x'_j\in A^-_a$.
    \end{lem}	 
	\begin{proof}
		Note that $y_i$ is linearly expressible in sum of two elements of the form $x+x^a$. The assertion of the lemma follows from Lemmas \ref{cd} and \ref{dc}.
		%Заметим, что $y_i$ разлагается в сумму двух элементов вида $x+x^a$. Из лемм \ref{cd} и \ref{dc} следует утверждение леммы.
	\end{proof} 
Thus, from Lemmas \ref{betia}, \ref{zbet} and \ref{alpbet} it follows that the products $x*y$ and $y*x$ lie in $A_a^- $, where $x\in A_a^+$ and $y_a\in A^-_a$.
%	 Таким образом из лемм \ref{betia}, \ref{zbet} и \ref{alpbet} следует,что произведения $x*y,y*x\in A_a^- $, где $x\in A_a^+$ и $y_a\in A^-$.
	 
	 We show that $A_a^+$ is a subalgebra. Note that $A_a^+$ has a basis consisting of $a$ and vectors of the form $l+l^a$, where $l\in T$. 
	\begin{lem}\label{pp}
	 	Let $x,y\in T$. Then $(x+x^a)*(y+y^a)\in A^+_a$.
	\end{lem} 
	\begin{proof}
We have $$(x+x^a)*(y+y^a)=x*y+x^a*y^a+x^a*y+x*y^a$$.
The assertion of the lemma follows from Lemma \ref{xxagen}.
%		Имеем $$(x+x^a)*(y+y^a)=x*y+x^a*y^a+x^a*y+x*y^a$$.
%	 	Из леммы \ref{xxagen} следует утверждение леммы.
	\end{proof}

\begin{lem}\label{pa}
	Let $x\in A^+_a$. We have $x*a\in A^+_a$.
\end{lem}
\begin{proof}
We can assume that $x$ lies in a $2$-generated subalgebra. The assertion of the lemma follows from Lemma \ref{prod1}.
%Можно считать, что $x$ лежит в $2$-порожденной подалгебре. Утверждение леммы следует из леммы \ref{prod1}. 
\end{proof}

Thus, it is proved that $(A^+_a,A^-_a)$ is a $\Z_2$-grading of the algebra $A$.
\end{proof}

\end{document}
